\paragraph{Pseudorandom generator (author Introduction).} Let $\mathcal F$ be a class of functions. A pseudorandom generator (PRG) for $\mathcal F$ is an algorithm $G:\{0,1\}^s\to\{0,1\}^n$ that maps a short random seed $x$ into a longer string $G(x)$ that appears random to every distinguisher in $\mathcal F$. More specifically, we say that $G$ $\epsilon$-fools $\mathcal F$, if for every $f\in\mathcal F$, it holds $|\mathbb E_{x\sim U_s}[f(G(x))]-\mathbb E_{x\sim U_n}[f(x)]|\le\epsilon$.
\section*{Author preliminaries (Section 3)}
We use the term ``restriction'' and ``partial assignment'' interchangeably. Both of them refer to a string of the form $\rho\in\{0,1,\star\}$. Here, if $\rho_i=0/1$, it means the $i$-th bit of $\rho$ is fixed to that value. Otherwise, the $i$-th bit of $\rho$ is unfixed. For two partial assignments $\rho,\sigma\in\{0,1,\star\}$, define their composition $\rho\circ\sigma$ as:
\[(\rho\circ\sigma)_i=\begin{cases}\rho_i&\rho_i\ne\star\\\sigma_i&\text{o.w.}\end{cases}.\]
Note that the left partial assignment always has a higher priority than the right one. Let $\Lambda\subseteq[n]$ be a set. For two partial assignments $\rho,\sigma\in\{0,1,\star\}^n$, let $\Lambda[\rho,\sigma]$ be the assignment defined as
\[\Lambda[\rho,\sigma]_i=\begin{cases}\rho_i&i\notin\Lambda\\\sigma_i&i\in\Lambda\end{cases}.\]
Let $f:\{0,1\}^n\to\{0,1\}$ be a function and $\rho\in\{0,1,\star\}^n$ be a partial assignment. We use $f|_\rho:\{0,1\}^n\to\{0,1\}$ to denote the restriction of $f$ on $\rho$. That is, $f|_\rho(y):=f(\rho\circ y)$.
$\mathrm{AC}^0_d[s(n)]$ denotes the family of Boolean circuits of size at most $s(n)$ and depth at most $d$. Such circuits can have AND, OR, NOT gates. Here, NOT gates do not count in the depth, and AND, OR gates can have unbounded fan-in. We measure the size of a circuit by the total number of wires (including input wires) in it. For the case of $d=2$, we also use the terms DNF and CNF to refer to $\mathrm{OR}\circ\mathrm{AND}$ and $\mathrm{AND}\circ\mathrm{OR}$ circuits respectively. The width of a DNF or CNF is defined as the maximum of its bottom fan-in. We also use $k$-DNF (resp. $k$-CNF) to denote DNF (resp. CNF) of width at most $k$.
A decision tree $T$ is a binary tree. Each inner node of $T$ is labelled with an index $i\in[n]$, and has exactly two children, which are labelled with 0 and 1. Each leaf of $T$ is labelled with a Boolean value $b\in\{0,1\}$. A decision tree $T$ computes a function in the following manner: on an input $x\in\{0,1\}^n$, we start from the root of $T$. In each turn we observe the index $i$ of current node, query $x_i$ and move to the left/right child depending on the bit $x_i$ we received. Once we reach a leaf with label $b$, we output $T(x)=b$. The depth of a decision tree is the length of the longest path from root to any leaf. We also consider a special decision tree model for a list of functions.
\paragraph{Definition 4.} Let $t,w,n,m\ge1$ be integers. Let $\mathcal F=\{F_1,\ldots,F_m\}$ be a list of functions mapping from $\{0,1\}^n$ to $\{0,1\}$. A $w$-partial depth-$t$ decision tree for $\mathcal F$ is a depth-$t$ decision tree $T$ satisfying the following. For every $F_i\in\mathcal F$ and every leaf $\ell$ of $T$, let $\alpha\in\{0,1,\star\}^n$ be the partial assignment that corresponds to $\ell$ (That is, $\alpha_i$ equals $\star$ if $T$ does not query $x_i$ before reaching $\ell$, otherwise $\alpha_i$ equals to the value that leads $T$ to move towards $\ell$). Then it holds that $\mathrm{DT}(F_i|_\alpha)\le w$.
\paragraph{Definition 5.} Let $n,m$ be two integers. Let $H$ be a distribution over hash functions mapping $\{0,1\}^n$ into $\{0,1\}^m$. We say that $H$ is $k$-wise independent, if for any $k$ input-output pairs $(x_1,y_1),\ldots,(x_k,y_k)\in\{0,1\}^n\times\{0,1\}^m$ where $x_1,\ldots,x_t$ are distinct, it holds that
\[\Pr_{h\sim H}[\forall i\in[k],h(x_i)=y_i]=2^{-km}.\]
We have the following standard construction of bounded independence hash functions (check e.g., [28, Chapter 3.5.5]).
\paragraph{Lemma 6.} For every $n,m,k\ge1$, there is an explicit $k$-wise independent hash functions $H$ that maps $\{0,1\}^n$ into $\{0,1\}^m$. One can sample a function in $H$ using $O(k(n+m))$ random bits.
We also consider a weaker notion of pseudorandomness called ``$k$-wise $p$-boundedness'', first defined and studied by [18].
\paragraph{Definition 7.} Suppose $\Lambda$ is a random subset of $[n]$. We say that $\Lambda$ is $k$-wise $p$-bounded if for any set $B\subseteq[n]$ of size at most $k$, it holds that $\Pr_\Lambda[B\subseteq\Lambda]\le p^{|B|}$. For intuition, if we sample $\Lambda$ by independently including each $i\in\Lambda$ with probability at most $p$, then $\Lambda$ is $n$-wise $p$-bounded.
\paragraph{Theorem 8 ([5,26]).} For every $m,n\in\mathbb N$ such that $m\ge n$ and every $\epsilon>0$, there is an $\epsilon$-PRG for $\mathrm{AC}^0_2[m]$ circuits (namely, CNF/DNF formulae) with seed length $O(\log(m)\log(m/\epsilon)\log\log(m))$.
